\sloppy

\begin{abstract}
    In multiethnic societies, interethnic marriage serves as a critical lens for
    understanding cultural integration and status hierarchies among ethnic
    groups.
    Using data from China's Census (1982-2010), we examine how these dynamics
    unfold in a context where traditional ethnic hierarchies intersect with
    state interventions through preferential policies.
    Although uncommon, intermarriage with the Han majority gradually increased,
    occurring most frequently among the Manchu, Mongolian, and Southern
    minorities, and least frequently among the Kazakh and Uyghur.
    After controlling for ethnic compositions, all minority ethnic groups
    exhibited preferences against intermarriage with the Han, with the strongest
    disinclination observed among groups maintaining distinct ethnoreligious
    identities (e.g., Kazakh, Uyghur).
    However, these cultural boundaries became increasingly permeable over time
    for most groups except the Manchu.
    Our analysis of educational patterns reveals two distinct mechanisms for
    boundary crossing: highly assimilated groups (e.g., Hui, Manchu, Mongolian)
    through educational assortative mating, others through status
    exchange--Koreans leveraged their educational advantage to marry Han
    spouses, while Han Chinese needed higher education to marry into Tibetan and
    Southern minorities.
    These findings reveal how marriage choices embody complex negotiations
    between cultural preservation, social mobility, and state-structured
    incentives, providing insights into the persistence and malleability of
    ethnic boundaries in contemporary societies.
\end{abstract}